% PyObj Z3 Theory + Semantic Algebra (Ch 25-26)
% Lines 5194-6844 of main.tex
% This file is for agent editing — will be merged back

\chapter{The PyObj Z3 Theory: A Complete Semantic Model of Python}
\label{ch:pyobj}

This chapter defines the Z3 theory that models Python's value space,
operations, and execution semantics.  The theory is designed to be
\emph{maximally expressive} (capturing as much of Python's semantics
as possible) while \emph{preserving decidability} (every previously
decidable query remains decidable after adding new features).

The key design principle: \textbf{new expressiveness is added via
new constructors and new shared functions, never via quantifiers}.
This keeps all fiber-restricted queries in the quantifier-free
fragment QF\_DT + QF\_LIA + QF\_S, which is decidable.

\section{The Value Universe}

\subsection{The Seven Constructors}

The universe of Python values is modeled as a Z3 algebraic data type
with seven constructors:

\begin{definition}[PyObj ADT]
\[
\PyObj ::=
  \mathsf{IntObj}(i : \Int) \mid
  \mathsf{BoolObj}(b : \Bool) \mid
  \mathsf{StrObj}(s : \mathsf{String}) \mid
  \mathsf{NoneObj} \mid
  \mathsf{Pair}(a : \PyObj, b : \PyObj) \mid
  \mathsf{Ref}(\mathsf{addr} : \Int) \mid
  \mathsf{Bottom}
\]
\end{definition}

\begin{lstlisting}[caption={PyObj ADT: complete Z3 implementation}]
from z3 import *

D = Datatype('PyObj')
D.declare('IntObj',  ('ival', IntSort()))
D.declare('BoolObj', ('bval', BoolSort()))
D.declare('StrObj',  ('sval', StringSort()))
D.declare('NoneObj')
D.declare('Pair', ('fst', D), ('snd', D))
D.declare('Ref',  ('addr', IntSort()))
D.declare('Bottom')
S = D.create()

# Constructor shortcuts
def mkint(n):  return S.IntObj(IntVal(n))
def mkbool(b): return S.BoolObj(BoolVal(b))
def mkstr(s):  return S.StrObj(StringVal(s))
def mklist(*elts):
    r = S.NoneObj
    for e in reversed(elts): r = S.Pair(e, r)
    return r
\end{lstlisting}

\begin{center}
\begin{tabular}{l l l l}
  \textbf{Constructor} & \textbf{Python type(s)} & \textbf{Z3 field} &
  \textbf{Exactness} \\
  \hline
  $\mathsf{IntObj}$ & \texttt{int} & Z3 IntSort (unbounded) & Exact \\
  $\mathsf{BoolObj}$ & \texttt{bool} & Z3 BoolSort & Exact \\
  $\mathsf{StrObj}$ & \texttt{str} & Z3 StringSort & Exact \\
  $\mathsf{NoneObj}$ & \texttt{None} & (no field) & Exact \\
  $\mathsf{Pair}$ & \texttt{list}, \texttt{tuple} & cons cell & Structural \\
  $\mathsf{Ref}$ & mutable objects & heap address & Abstract \\
  $\mathsf{Bottom}$ & exceptions, $\bot$ & (no field) & Abstract \\
\end{tabular}
\end{center}

\subsection{Why Seven Constructors Suffice}

Python has hundreds of types (\texttt{int}, \texttt{float},
\texttt{complex}, \texttt{str}, \texttt{bytes}, \texttt{list},
\texttt{tuple}, \texttt{set}, \texttt{frozenset}, \texttt{dict},
\texttt{deque}, \texttt{defaultdict}, \texttt{Counter},
\texttt{OrderedDict}, and every user-defined class).  Our seven
constructors cover them via \emph{abstraction}:

\begin{center}
\small
\begin{tabular}{l l l}
  \textbf{Python type} & \textbf{PyObj constructor} &
  \textbf{Justification} \\
  \hline
  \texttt{int} & $\mathsf{IntObj}$ & Exact: Z3 IntSort is
    unbounded, matching Python's arbitrary-precision ints \\
  \texttt{float} & $\mathsf{IntObj}$ & Approximate: we model
    float as int (sound for integer-valued floats) \\
  \texttt{complex} & $\mathsf{Pair}(\mathsf{IntObj},
    \mathsf{IntObj})$ & Real and imaginary parts \\
  \texttt{bool} & $\mathsf{BoolObj}$ & Exact \\
  \texttt{str} & $\mathsf{StrObj}$ & Exact: Z3 StringSort
    models Unicode strings with concat, length, etc. \\
  \texttt{bytes} & $\mathsf{StrObj}$ & Approximate \\
  \texttt{None} & $\mathsf{NoneObj}$ & Exact (singleton) \\
  \texttt{list} & $\mathsf{Pair}$-chain & Structural: $[1,2,3]$
    = Pair(1, Pair(2, Pair(3, NoneObj))) \\
  \texttt{tuple} & $\mathsf{Pair}$-chain & Same encoding as
    list (Python tuples and lists have same structure) \\
  \texttt{set} & $\mathsf{Ref}$ & Abstract: modeled by identity
    (heap address), operations via shared functions \\
  \texttt{frozenset} & $\mathsf{Ref}$ & Abstract \\
  \texttt{dict} & $\mathsf{Ref}$ & Abstract \\
  \texttt{deque} & $\mathsf{Ref}$ & Abstract \\
  \texttt{defaultdict} & $\mathsf{Ref}$ & Abstract \\
  \texttt{Counter} & $\mathsf{Ref}$ & Abstract \\
  \texttt{OrderedDict} & $\mathsf{Ref}$ & Abstract \\
  user class instance & $\mathsf{Ref}$ & Abstract \\
  function/lambda & $\mathsf{Ref}$ & Abstract (callable ref) \\
  generator & $\mathsf{Ref}$ & Abstract (stateful iterator) \\
  exception & $\mathsf{Bottom}$ & Abstract (control flow) \\
  \texttt{NotImplemented} & $\mathsf{Bottom}$ & Abstract \\
\end{tabular}
\end{center}

The abstraction is \emph{sound}: if the Z3 theory proves two
programs equivalent, they are genuinely equivalent for all concrete
Python values that the constructors faithfully represent.  The
abstraction may be \emph{incomplete}: programs that differ only on
types mapped to $\mathsf{Ref}$ (e.g., two different dict
implementations) may not be distinguished.

\subsection{Decidability Preservation}

\begin{proposition}[Adding Constructors Preserves Decidability]
\label{prop:adt-decidable}
The theory QF\_DT (quantifier-free datatypes) is decidable for any
finite set of constructors.  Adding a new constructor to the $\PyObj$
ADT does not break decidability --- it adds new cases to pattern
matching but keeps the theory in QF\_DT.
\end{proposition}

This means we can extend $\PyObj$ with new constructors (e.g.,
$\mathsf{FloatObj}(\mathsf{real} : \mathsf{RealSort})$,
$\mathsf{SetObj}(\mathsf{elems} : \mathsf{ArraySort})$) without
breaking decidability of existing queries.  We choose seven
constructors for practical efficiency (fewer cases in RecFunctions),
not for theoretical necessity.

\section{The Type Fibration}

\subsection{TypeTag Sort}

The type tag function is the fibration structure map:

\begin{definition}[Type Fibration]
$\mathsf{tag} : \PyObj \to \mathsf{TypeTag}$ where
$\mathsf{TypeTag} = \{\mathsf{TInt}, \mathsf{TBool}, \mathsf{TStr},
\mathsf{TNone}, \mathsf{TPair}, \mathsf{TRef}, \mathsf{TBot}\}$
is defined by:
\[
  \mathsf{tag}(v) = \begin{cases}
    \mathsf{TInt} & \text{if } v = \mathsf{IntObj}(\ldots) \\
    \mathsf{TBool} & \text{if } v = \mathsf{BoolObj}(\ldots) \\
    \mathsf{TStr} & \text{if } v = \mathsf{StrObj}(\ldots) \\
    \mathsf{TNone} & \text{if } v = \mathsf{NoneObj} \\
    \mathsf{TPair} & \text{if } v = \mathsf{Pair}(\ldots) \\
    \mathsf{TRef} & \text{if } v = \mathsf{Ref}(\ldots) \\
    \mathsf{TBot} & \text{if } v = \mathsf{Bottom}
  \end{cases}
\]
\end{definition}

This is implemented as a Z3 RecFunction and a Z3 EnumSort.  The
RecFunction dispatches on the constructor, returning the corresponding
enum constant.

\begin{lstlisting}[caption={TypeTag: complete Z3 implementation}]
TTag, (TInt, TBool, TStr, TNone, TPair, TRef, TBot) = \
    EnumSort('TTag', ['TInt','TBool','TStr',
                      'TNone','TPair','TRef','TBot'])

tag = RecFunction('tag', S, TTag)
x = Const('_tag_x', S)
RecAddDefinition(tag, [x],
    If(S.is_IntObj(x),  TInt,
    If(S.is_BoolObj(x), TBool,
    If(S.is_StrObj(x),  TStr,
    If(S.is_NoneObj(x), TNone,
    If(S.is_Pair(x),    TPair,
    If(S.is_Ref(x),     TRef,
                        TBot)))))))
\end{lstlisting}

\subsection{Fiber Projection Lemmas}

The following lemmas are used repeatedly in the \v{C}ech complex
construction.  They are all decidable consequences of the ADT theory.

\begin{lemma}[Fiber Determines Constructor]
\label{lem:fiber-constructor}
$\mathsf{tag}(v) = \mathsf{TInt} \iff \exists n : \Int.\;
v = \mathsf{IntObj}(n)$.

More generally, knowing the tag determines the constructor shape,
which Z3 exploits via the $\mathsf{is\_IntObj}$,
$\mathsf{is\_BoolObj}$, \ldots tester functions.
\end{lemma}

\begin{lemma}[Fiber Disjointness]
\label{lem:fiber-disjoint}
The fibers are pairwise disjoint: $\mathsf{tag}(v) = \tau_1 \land
\mathsf{tag}(v) = \tau_2 \implies \tau_1 = \tau_2$.
This is automatic from the ADT theory (each value has exactly one
constructor).
\end{lemma}

\begin{lemma}[Fiber Exhaustion]
The fibers cover all values: $\bigvee_{\tau \in \mathsf{TypeTag}}
\mathsf{tag}(v) = \tau$.  This is the ``no junk'' property of
inductive types.
\end{lemma}


\section{Operations: Full Python Semantics in Z3}

\subsection{Binary Operations}

The $\mathsf{binop}$ RecFunction encodes Python's binary operations
with full type dispatch:

\begin{definition}[Binary Operation Semantics]
$\mathsf{binop} : \Int \times \PyObj \times \PyObj \to \PyObj$
where the first argument is an operation tag.

\textbf{Integer $\times$ Integer} ($\mathsf{tag}(a) = \mathsf{tag}(b)
= \mathsf{TInt}$):
\begin{align}
  \mathsf{binop}(\mathsf{ADD}, \mathsf{IntObj}(m),
    \mathsf{IntObj}(n)) &= \mathsf{IntObj}(m + n) \\
  \mathsf{binop}(\mathsf{SUB}, \mathsf{IntObj}(m),
    \mathsf{IntObj}(n)) &= \mathsf{IntObj}(m - n) \\
  \mathsf{binop}(\mathsf{MUL}, \mathsf{IntObj}(m),
    \mathsf{IntObj}(n)) &= \mathsf{IntObj}(m \times n) \\
  \mathsf{binop}(\mathsf{FLOORDIV}, \mathsf{IntObj}(m),
    \mathsf{IntObj}(n)) &= \begin{cases}
      \mathsf{IntObj}(m \div n) & n \neq 0 \\
      \mathsf{Bottom} & n = 0
    \end{cases} \\
  \mathsf{binop}(\mathsf{MOD}, \mathsf{IntObj}(m),
    \mathsf{IntObj}(n)) &= \begin{cases}
      \mathsf{IntObj}(m \bmod n) & n \neq 0 \\
      \mathsf{Bottom} & n = 0
    \end{cases} \\
  \mathsf{binop}(\mathsf{LT}, \mathsf{IntObj}(m),
    \mathsf{IntObj}(n)) &= \mathsf{BoolObj}(m < n) \\
  \mathsf{binop}(\mathsf{LE}, \ldots) &= \mathsf{BoolObj}(m \leq n) \\
  \mathsf{binop}(\mathsf{GT}, \ldots) &= \mathsf{BoolObj}(m > n) \\
  \mathsf{binop}(\mathsf{GE}, \ldots) &= \mathsf{BoolObj}(m \geq n) \\
  \mathsf{binop}(\mathsf{EQ}, \ldots) &= \mathsf{BoolObj}(m = n) \\
  \mathsf{binop}(\mathsf{NE}, \ldots) &= \mathsf{BoolObj}(m \neq n) \\
  \mathsf{binop}(\mathsf{MAX}, \ldots) &= \mathsf{If}(m \geq n,
    \mathsf{IntObj}(m), \mathsf{IntObj}(n)) \\
  \mathsf{binop}(\mathsf{MIN}, \ldots) &= \mathsf{If}(m \leq n,
    \mathsf{IntObj}(m), \mathsf{IntObj}(n))
\end{align}

\textbf{String $\times$ String} ($\mathsf{tag}(a) = \mathsf{tag}(b)
= \mathsf{TStr}$):
\begin{align}
  \mathsf{binop}(\mathsf{ADD}, \mathsf{StrObj}(s),
    \mathsf{StrObj}(t)) &= \mathsf{StrObj}(\mathsf{Concat}(s, t)) \\
  \mathsf{binop}(\mathsf{EQ}, \mathsf{StrObj}(s),
    \mathsf{StrObj}(t)) &= \mathsf{BoolObj}(s = t) \\
  \mathsf{binop}(\mathsf{NE}, \ldots) &= \mathsf{BoolObj}(s \neq t)
\end{align}

\textbf{Cross-type equality}: For any $a, b$ with
$\mathsf{tag}(a) \neq \mathsf{tag}(b)$:
\begin{align}
  \mathsf{binop}(\mathsf{EQ}, a, b) &= \mathsf{BoolObj}(a = b)
    && \text{(structural equality)} \\
  \mathsf{binop}(\mathsf{NE}, a, b) &= \mathsf{BoolObj}(a \neq b)
\end{align}

\textbf{All other cross-type combinations}: $\mathsf{Bottom}$
(TypeError in Python).
\end{definition}

\begin{proposition}[binop Decidability]
\label{prop:binop-decidable}
On the $\mathsf{TInt}$ fiber, $\mathsf{binop}$ reduces to Z3's
native $\mathsf{IntSort}$ operations (decidable QF\_LIA for linear
operations, semi-decidable QF\_NIA for multiplication).
On the $\mathsf{TStr}$ fiber, it reduces to Z3's string theory
(decidable QF\_S for concat and equality).
On other fibers, it returns $\mathsf{Bottom}$ (decidable constant).

Therefore: $\mathsf{binop}$ does not break decidability on any
fiber.
\end{proposition}

\begin{lstlisting}[caption={binop: complete Z3 implementation}]
ADD=1; SUB=2; MUL=3; FLOORDIV=5; MOD=6
LT=20; LE=21; GT=22; GE=23; EQ=24; NE=25
MAX=30; MIN=31

binop = RecFunction('binop', IntSort(), S, S, S)
op = Int('_bo')
a, b = Const('_ba', S), Const('_bb', S)
ai, bi = S.ival(a), S.ival(b)
ii = And(S.is_IntObj(a), S.is_IntObj(b))
ss = And(S.is_StrObj(a), S.is_StrObj(b))

int_result = \
  If(op==ADD, S.IntObj(ai+bi),
  If(op==SUB, S.IntObj(ai-bi),
  If(op==MUL, S.IntObj(ai*bi),
  If(op==FLOORDIV, If(bi!=0, S.IntObj(ai/bi), S.Bottom),
  If(op==MOD, If(bi!=0, S.IntObj(ai%bi), S.Bottom),
  If(op==LT, S.BoolObj(ai<bi),
  If(op==LE, S.BoolObj(ai<=bi),
  If(op==GT, S.BoolObj(ai>bi),
  If(op==GE, S.BoolObj(ai>=bi),
  If(op==EQ, S.BoolObj(ai==bi),
  If(op==NE, S.BoolObj(ai!=bi),
  If(op==MAX, If(ai>=bi, a, b),
  If(op==MIN, If(ai<=bi, a, b),
  S.Bottom)))))))))))))

str_result = \
  If(op==ADD, S.StrObj(Concat(S.sval(a), S.sval(b))),
  If(op==EQ, S.BoolObj(S.sval(a)==S.sval(b)),
  If(op==NE, S.BoolObj(S.sval(a)!=S.sval(b)),
  S.Bottom)))

RecAddDefinition(binop, [op, a, b],
    If(ii, int_result,
    If(ss, str_result,
    If(op==EQ, S.BoolObj(a==b),
    If(op==NE, S.BoolObj(a!=b),
    S.Bottom)))))
\end{lstlisting}

\subsection{Unary Operations}

\begin{definition}[Unary Operation Semantics]
$\mathsf{unop} : \Int \times \PyObj \to \PyObj$

\textbf{On IntObj}:
\begin{align}
  \mathsf{unop}(\mathsf{NEG}, \mathsf{IntObj}(n)) &=
    \mathsf{IntObj}(-n) \\
  \mathsf{unop}(\mathsf{ABS}, \mathsf{IntObj}(n)) &=
    \mathsf{IntObj}(\mathsf{If}(n \geq 0, n, -n)) \\
  \mathsf{unop}(\mathsf{NOT}, \mathsf{IntObj}(n)) &=
    \mathsf{BoolObj}(n = 0) \\
  \mathsf{unop}(\mathsf{BOOL}, \mathsf{IntObj}(n)) &=
    \mathsf{BoolObj}(n \neq 0) \\
  \mathsf{unop}(\mathsf{INT}, \mathsf{IntObj}(n)) &=
    \mathsf{IntObj}(n)
\end{align}

\textbf{On BoolObj}:
\begin{align}
  \mathsf{unop}(\mathsf{NOT}, \mathsf{BoolObj}(b)) &=
    \mathsf{BoolObj}(\lnot b) \\
  \mathsf{unop}(\mathsf{BOOL}, \mathsf{BoolObj}(b)) &=
    \mathsf{BoolObj}(b) \\
  \mathsf{unop}(\mathsf{INT}, \mathsf{BoolObj}(b)) &=
    \mathsf{IntObj}(\mathsf{If}(b, 1, 0))
\end{align}

\textbf{On StrObj}:
\begin{align}
  \mathsf{unop}(\mathsf{LEN}, \mathsf{StrObj}(s)) &=
    \mathsf{IntObj}(\mathsf{Length}(s)) \\
  \mathsf{unop}(\mathsf{BOOL}, \mathsf{StrObj}(s)) &=
    \mathsf{BoolObj}(\mathsf{Length}(s) > 0)
\end{align}

\textbf{On NoneObj}: $\mathsf{unop}(\mathsf{BOOL}, \mathsf{NoneObj})
= \mathsf{BoolObj}(\bot)$; $\mathsf{unop}(\mathsf{NOT},
\mathsf{NoneObj}) = \mathsf{BoolObj}(\top)$.

\textbf{On Pair}: $\mathsf{unop}(\mathsf{BOOL},
\mathsf{Pair}(a, b)) = \mathsf{BoolObj}(\top)$ (non-empty list is
truthy).
\end{definition}

\subsection{The Truthiness Protocol}

Python's $\mathsf{\_\_bool\_\_}$ protocol is central to control flow.
Every value has a truth value:

\begin{definition}[Truthiness]
$\mathsf{truthy} : \PyObj \to \Bool$
\begin{align}
  \mathsf{truthy}(\mathsf{IntObj}(n)) &= (n \neq 0) \\
  \mathsf{truthy}(\mathsf{BoolObj}(b)) &= b \\
  \mathsf{truthy}(\mathsf{StrObj}(s)) &= (\mathsf{Length}(s) > 0) \\
  \mathsf{truthy}(\mathsf{NoneObj}) &= \bot \\
  \mathsf{truthy}(\mathsf{Pair}(a, b)) &= \top \\
  \mathsf{truthy}(\mathsf{Bottom}) &= \bot \\
  \mathsf{truthy}(\mathsf{Ref}(\_)) &= \top
\end{align}
\end{definition}

\begin{remark}[Faithfulness of Truthiness]
This exactly matches CPython's \texttt{\_\_bool\_\_} for all modeled
types.  The only approximation is $\mathsf{Ref}$: we model all
reference objects as truthy, which is correct for non-empty containers
but wrong for empty containers stored as Refs.  This is sound for
the \v{C}ech analysis because the Ref fiber is handled abstractly
(via shared function symbols, not by evaluating truthiness).
\end{remark}

\begin{lstlisting}[caption={truthy: complete Z3 implementation}]
truthy = RecFunction('truthy', S, BoolSort())
x = Const('_tr', S)
RecAddDefinition(truthy, [x],
    If(S.is_IntObj(x),  S.ival(x) != 0,
    If(S.is_BoolObj(x), S.bval(x),
    If(S.is_StrObj(x),  Length(S.sval(x)) > 0,
    If(S.is_NoneObj(x), False,
    If(S.is_Bottom(x),  False,
    True))))))  # Pair, Ref are truthy
\end{lstlisting}

\subsection{The Canonical Fold (Nat-Recursor)}

\begin{definition}[Canonical Fold]
$\fold : \Int \times \Int \times \Int \times \PyObj \to \PyObj$
\begin{align}
  \fold(\op, i, \mathsf{stop}, \mathsf{acc}) &= \begin{cases}
    \mathsf{acc} & \text{if } i \geq \mathsf{stop} \\
    \mathsf{binop}(\op, \mathsf{IntObj}(i),
      \fold(\op, i+1, \mathsf{stop}, \mathsf{acc}))
      & \text{otherwise}
  \end{cases}
\end{align}
\end{definition}

\begin{lstlisting}[caption={fold: complete Z3 implementation}]
fold = RecFunction('fold',
    IntSort(), IntSort(), IntSort(), S, S)
op, i, stop = Ints('_fo _fi _fs')
acc = Const('_fa', S)
RecAddDefinition(fold, [op, i, stop, acc],
    If(i >= stop, acc,
       binop(op, S.IntObj(i),
             fold(op, i+1, stop, acc))))
\end{lstlisting}

The fold is the \emph{universal} iteration primitive.  Both recursive
and iterative programs that operate by applying a binary operation
across a range compile to this single RecFunction, enabling
equivalence by syntactic identity (D1 path).

\begin{proposition}[Fold Decidability]
\label{prop:fold-decidable}
The fold does not break decidability of previously decidable queries.
When both programs compile to the \emph{same} fold term, the
comparison is syntactic equality (decidable, $O(1)$).  When they
compile to \emph{different} fold terms, Z3 handles the comparison
via its RecFunction reasoning (semi-decidable in general, but
decidable for ground arguments).
\end{proposition}

\subsection{Item Access (GETITEM)}

\begin{definition}[Subscript Semantics]
$\mathsf{binop}(\mathsf{GETITEM}, \mathsf{collection}, \mathsf{index})$
implements Python's $\mathsf{\_\_getitem\_\_}$:

\textbf{On Pair-chains} (list indexing):
\begin{align}
  \mathsf{binop}(\mathsf{GETITEM}, \mathsf{Pair}(a, d),
    \mathsf{IntObj}(0)) &= a \\
  \mathsf{binop}(\mathsf{GETITEM}, \mathsf{Pair}(a, d),
    \mathsf{IntObj}(n)) &= \mathsf{binop}(\mathsf{GETITEM}, d,
    \mathsf{IntObj}(n-1)) \quad n > 0
\end{align}

\textbf{On StrObj} (character indexing):
\[
  \mathsf{binop}(\mathsf{GETITEM}, \mathsf{StrObj}(s),
    \mathsf{IntObj}(n)) = \mathsf{StrObj}(\mathsf{SubString}(s, n, 1))
\]

\textbf{On Ref} (dict/set access): uses shared function symbol
$\mathsf{py\_meth\_\_\_getitem\_\_}$.
\end{definition}


\section{Container Semantics via Shared Functions}

For types modeled by $\mathsf{Ref}$ (dicts, sets, deques, etc.),
operations are modeled via shared function symbols with
\emph{algebraic axioms}.  This is more expressive than the ADT
constructors (we capture behavioral properties without modeling
internal structure) and preserves decidability (shared functions
are uninterpreted; axioms are optional path constructors).

\subsection{Dictionary Operations}

\begin{definition}[Dict Semantic Functions]
\begin{align}
  \mathsf{py\_call\_dict} &: \PyObj \to \PyObj
    && \text{construct from pairs} \\
  \mathsf{py\_meth\_\_\_getitem\_\_} &: \PyObj \times \PyObj \to \PyObj
    && \text{key lookup} \\
  \mathsf{py\_meth\_\_\_setitem\_\_} &: \PyObj \times \PyObj \times
    \PyObj \to \PyObj
    && \text{key assignment (returns new dict)} \\
  \mathsf{py\_meth\_get} &: \PyObj \times \PyObj \to \PyObj
    && \text{get with default} \\
  \mathsf{py\_meth\_keys} &: \PyObj \to \PyObj
    && \text{key set} \\
  \mathsf{py\_meth\_values} &: \PyObj \to \PyObj
    && \text{value collection} \\
  \mathsf{py\_meth\_items} &: \PyObj \to \PyObj
    && \text{key-value pairs} \\
  \mathsf{py\_meth\_update} &: \PyObj \times \PyObj \to \PyObj
    && \text{merge dicts} \\
  \mathsf{py\_meth\_pop} &: \PyObj \times \PyObj \to \PyObj
    && \text{remove key}
\end{align}
\end{definition}

These are \emph{uninterpreted} by default (Z3 knows nothing about
them except that they are functions).  Axioms can be added as path
constructors:

\begin{align}
  &\forall d, k, v.\; \mathsf{py\_meth\_\_\_getitem\_\_}(
    \mathsf{py\_meth\_\_\_setitem\_\_}(d, k, v), k) = v
    && \text{(set-then-get)} \\
  &\forall d, k_1, k_2, v.\; k_1 \neq k_2 \implies \nonumber \\
  &\quad \mathsf{py\_meth\_\_\_getitem\_\_}(
    \mathsf{py\_meth\_\_\_setitem\_\_}(d, k_1, v), k_2) =
    \mathsf{py\_meth\_\_\_getitem\_\_}(d, k_2)
    && \text{(set-then-get-other)}
\end{align}

\begin{proposition}[Dict Axioms Preserve Decidability]
These axioms are universally quantified over $\PyObj$, so they
technically exit the QF fragment.  However:
\begin{enumerate}
  \item On a specific fiber, the quantifiers are bounded (the
    variables have fixed tags), so E-matching terminates.
  \item The axioms are \emph{conditional} ($k_1 \neq k_2$ guard),
    which limits instantiation.
  \item They are added only when both programs actually use dict
    operations (axiom relevance filtering).
\end{enumerate}
In practice, these axioms do not cause Z3 timeouts on the benchmark.
\end{proposition}

\subsection{List Operations}

\begin{definition}[List Semantic Functions]
\begin{align}
  \mathsf{py\_mut\_append} &: \PyObj \times \PyObj \to \PyObj
    && \text{append element (mutation)} \\
  \mathsf{py\_mut\_extend} &: \PyObj \times \PyObj \to \PyObj
    && \text{extend with another list} \\
  \mathsf{py\_mut\_insert} &: \PyObj \times \PyObj \times \PyObj
    \to \PyObj
    && \text{insert at index} \\
  \mathsf{py\_mut\_pop} &: \PyObj \to \PyObj
    && \text{remove last element} \\
  \mathsf{py\_mut\_sort} &: \PyObj \to \PyObj
    && \text{in-place sort (mutation!)} \\
  \mathsf{py\_mut\_reverse} &: \PyObj \to \PyObj
    && \text{in-place reverse (mutation!)}
\end{align}
\end{definition}

Note the distinction between \emph{mutation} operations (prefixed
$\mathsf{mut\_}$, which modify a Ref in place) and \emph{pure}
operations (like $\mathsf{py\_sorted}$, which return a new value).
This distinction is critical for the effect discipline
(Chapter~\ref{ch:effects}): mutation operations use different shared
symbols than their pure counterparts, so Z3 correctly distinguishes
$\mathsf{lst.sort()}$ (mutation, returns None) from
$\mathsf{sorted(lst)}$ (pure, returns new list).

\subsection{Set Operations}

\begin{definition}[Set Semantic Functions]
\begin{align}
  \mathsf{py\_mut\_add} &: \PyObj \times \PyObj \to \PyObj
    && \text{add element} \\
  \mathsf{py\_mut\_discard} &: \PyObj \times \PyObj \to \PyObj
    && \text{remove element} \\
  \mathsf{py\_call\_union} &: \PyObj \times \PyObj \to \PyObj
    && \text{set union} \\
  \mathsf{py\_call\_intersection} &: \PyObj \times \PyObj \to \PyObj
    && \text{set intersection} \\
  \mathsf{py\_call\_difference} &: \PyObj \times \PyObj \to \PyObj
    && \text{set difference}
\end{align}
\end{definition}

Set axioms (optional path constructors):
\begin{align}
  &\mathsf{py\_call\_union}(a, b) = \mathsf{py\_call\_union}(b, a)
    && \text{(union commutativity)} \\
  &\mathsf{py\_call\_intersection}(a, b) =
    \mathsf{py\_call\_intersection}(b, a)
    && \text{(intersection commutativity)} \\
  &\mathsf{py\_call\_union}(a, a) = a
    && \text{(union idempotence)}
\end{align}


\section{The Heap Model}

Mutable Python objects live on a \emph{heap} --- a global store that
maps addresses to values.  Two references to the same address are
\emph{aliases}: modifying one is visible through the other.

\subsection{Heap Representation}

\begin{definition}[Abstract Heap]
The heap is modeled \emph{implicitly} through the $\mathsf{Ref}$
constructor and mutation operations:
\begin{itemize}
  \item $\mathsf{Ref}(\mathsf{addr})$ represents a heap-allocated
    object at address $\mathsf{addr}$.
  \item Fresh allocations use fresh addresses:
    $\mathsf{mkref}() = \mathsf{Ref}(\mathsf{next\_uid})$
    where $\mathsf{next\_uid}$ is a monotonically increasing counter.
  \item Mutation operations ($\mathsf{mut\_append}$,
    $\mathsf{mut\_sort}$, etc.) return a \emph{new version} of the
    object: $\mathsf{mut\_op}(\mathsf{Ref}(a), \ldots) = v'$
    where $v'$ represents the modified object.
\end{itemize}
\end{definition}

\begin{remark}[SSA for the Heap]
This is essentially \emph{Static Single Assignment} (SSA) form for
heap operations: each mutation produces a new value, and the
environment tracks the latest version.  Aliasing is modeled by
having multiple environment entries point to the same $\mathsf{Ref}$.
\end{remark}

\subsection{Aliasing and the Effect Discipline}

\begin{definition}[Aliasing Semantics]
Two variables $x$ and $y$ are \emph{aliases} if
$\mathsf{env}[x] = \mathsf{env}[y] = \mathsf{Ref}(a)$ for the same
address $a$.  A mutation on $x$ modifies the Ref, so $y$ (which
points to the same Ref) sees the modification.
\end{definition}

This is where the effect discipline matters:
\begin{itemize}
  \item $\mathsf{mut\_iadd}(\mathsf{Ref}(a), v)$ \emph{modifies}
    the Ref --- all aliases see the change.
  \item $\mathsf{py\_call\_\_\_add\_\_}(\mathsf{Ref}(a), v)$
    \emph{creates a new value} --- aliases are unaffected.
\end{itemize}

These are different shared symbols, so Z3 treats them as unrelated
functions.  This is how the framework correctly distinguishes
\texttt{lst += [1]} (mutates alias) from \texttt{lst = lst + [1]}
(rebinds variable).


\section{Python Control Flow in Z3}

\subsection{Conditional Expressions}

Python's \texttt{if}/\texttt{else} compiles to Z3's $\mathsf{If}$:
\[
  \den{\texttt{t if c else f}} = \mathsf{If}(\mathsf{truthy}(\den{c}),
    \den{t}, \den{f})
\]

This is a Z3 builtin (part of the core theory, always decidable).

\subsection{Short-Circuit Boolean Operators}

Python's \texttt{and}/\texttt{or} are short-circuit and return
\emph{values}, not booleans:
\begin{align}
  \den{\texttt{a and b}} &= \mathsf{If}(\mathsf{truthy}(\den{a}),
    \den{b}, \den{a}) \\
  \den{\texttt{a or b}} &= \mathsf{If}(\mathsf{truthy}(\den{a}),
    \den{a}, \den{b})
\end{align}

This faithfully models Python's semantics: \texttt{0 or "hello"}
returns \texttt{"hello"} (not \texttt{True}), and
\texttt{[] and 42} returns \texttt{[]} (not \texttt{False}).

\subsection{Exception Semantics via Bottom}

Exceptions are modeled by $\mathsf{Bottom}$:
\begin{itemize}
  \item Division by zero: $\mathsf{binop}(\mathsf{FLOORDIV}, x,
    \mathsf{IntObj}(0)) = \mathsf{Bottom}$.
  \item Type errors: $\mathsf{binop}(\mathsf{ADD},
    \mathsf{IntObj}(n), \mathsf{StrObj}(s)) = \mathsf{Bottom}$.
  \item Explicit raise: compiled to $\mathsf{Bottom}$.
\end{itemize}

$\mathsf{Bottom}$ propagates strictly:
$\mathsf{binop}(\op, \mathsf{Bottom}, x) = \mathsf{Bottom}$ and
$\mathsf{binop}(\op, x, \mathsf{Bottom}) = \mathsf{Bottom}$.

\texttt{try/except} blocks are modeled by wrapping the body in a
shared ``handler context'' function:
\[
  \den{\texttt{try: body except E: handler}} =
  \mathsf{py\_try\_h}(\den{\mathsf{body}})
\]
where $h$ is a hash of the handler signature.  The handler context
is a shared function, so two programs with the same try/except
structure get the same Z3 symbol.

\subsection{isinstance as Fiber Projection}

\begin{definition}[isinstance Semantics]
\[
  \den{\texttt{isinstance(x, T)}} = \mathsf{BoolObj}(\mathsf{tag}(\den{x}) = \tau_T)
\]
where $\tau_T$ is the tag corresponding to type $T$:

\begin{center}
\begin{tabular}{l l l l l l}
  \texttt{int} & $\to \mathsf{TInt}$ &
  \texttt{str} & $\to \mathsf{TStr}$ &
  \texttt{bool} & $\to \mathsf{TBool}$ \\
  \texttt{list} & $\to \mathsf{TPair}$ &
  \texttt{tuple} & $\to \mathsf{TPair}$ &
  \texttt{dict} & $\to \mathsf{TRef}$ \\
  \texttt{set} & $\to \mathsf{TRef}$ &
  \texttt{float} & $\to \mathsf{TInt}$ &
  \texttt{type} & $\to \mathsf{TRef}$
\end{tabular}
\end{center}

For tuple arguments (\texttt{isinstance(x, (int, str))}):
\[
  \den{\texttt{isinstance(x, (int, str))}} =
    \mathsf{BoolObj}(\mathsf{tag}(\den{x}) = \mathsf{TInt} \lor
    \mathsf{tag}(\den{x}) = \mathsf{TStr})
\]
\end{definition}

This is the \emph{fiber projection}: \texttt{isinstance} asks
which fiber of the type fibration the value lives in.  It is
purely decidable (tag comparison in EnumSort).

\subsection{Comprehensions}

List, set, and generator comprehensions compile to shared function
symbols keyed by a canonical hash of the element expression:

\begin{definition}[Comprehension Semantics]
\[
  \den{[\mathsf{elt}(x) \texttt{ for } x \texttt{ in } \mathsf{xs}]}
  = \mathsf{py\_comp}_{h(\mathsf{elt})}(\den{\mathsf{xs}})
\]
where $h(\mathsf{elt}) = \mathsf{hash}(\mathsf{ast.dump}(\mathsf{elt}))
\bmod 2^{31}$.
\end{definition}

Two comprehensions with the same element expression AST get the same
shared function symbol, regardless of whether they appear as a
comprehension or an explicit loop with append (D4 path).

\subsection{Attribute Access}

\begin{definition}[Attribute Semantics]
$\den{x.\mathsf{attr}} = \mathsf{py\_attr\_attr}(\den{x})$

This is a shared uninterpreted function.  No axioms are added by
default (we don't model Python's attribute lookup protocol in full).
But when both programs access the same attribute on the same object,
they get the same Z3 term (by sharing).
\end{definition}


\section{Shared Function Symbols (Univalence)}

\subsection{The Registry}

The \emph{shared function registry} ensures both programs use the
same Z3 function for the same Python operation.

\begin{definition}[Shared Function Registry]
$\mathcal{R} : (\mathsf{name} \times \mathsf{arity}) \to
\mathsf{Z3FuncDecl}$ maps operation names to canonical Z3 function
declarations.  The function $\mathsf{shared\_fn}(\mathsf{name},
\mathsf{arity})$ returns the canonical Z3 function, creating it on
first access.
\end{definition}

\subsection{Builtin Classification}

\begin{center}
\small
\begin{tabular}{l l l}
  \textbf{Category} & \textbf{Functions} & \textbf{Shared symbol} \\
  \hline
  Sorting & \texttt{sorted}, \texttt{reversed} &
    $\mathsf{py\_sorted}$, $\mathsf{py\_reversed}$ \\
  Collection & \texttt{list}, \texttt{tuple}, \texttt{set},
    \texttt{frozenset}, \texttt{dict} &
    $\mathsf{py\_list}$, etc. \\
  Aggregation & \texttt{sum}, \texttt{any}, \texttt{all},
    \texttt{min}, \texttt{max} &
    $\mathsf{py\_sum}$, etc. \\
  Higher-order & \texttt{map}, \texttt{filter}, \texttt{zip},
    \texttt{enumerate} &
    $\mathsf{py\_map}$, etc. \\
  Type conversion & \texttt{int}, \texttt{str}, \texttt{bool},
    \texttt{float} &
    via $\mathsf{unop}$ \\
  Type checking & \texttt{isinstance}, \texttt{type} &
    via $\mathsf{tag}$ \\
  String methods & \texttt{.upper}, \texttt{.lower},
    \texttt{.strip}, \texttt{.split}, \texttt{.join},
    \texttt{.replace}, \texttt{.find} &
    $\mathsf{py\_meth\_upper}$, etc. \\
  Dict methods & \texttt{.get}, \texttt{.keys}, \texttt{.values},
    \texttt{.items}, \texttt{.update}, \texttt{.pop} &
    $\mathsf{py\_meth\_get}$, etc. \\
  List methods & \texttt{.append}, \texttt{.extend}, \texttt{.sort},
    \texttt{.reverse}, \texttt{.pop}, \texttt{.insert} &
    $\mathsf{py\_mut\_append}$, etc. \\
  Mutation ops & \texttt{+=}, \texttt{*=}, etc.\ on containers &
    $\mathsf{py\_mut\_iadd}$, etc. \\
\end{tabular}
\end{center}

\subsection{Why Sharing Preserves Decidability}

\begin{proposition}[Shared Symbols Do Not Affect Decidability]
Replacing unique function symbols ($\mathsf{sorted}_1$,
$\mathsf{sorted}_5$) with shared symbols ($\mathsf{py\_sorted}$)
does not change the decidability status of any query.

In the QF fragment, all function symbols are uninterpreted --- Z3
treats them as ``black boxes'' regardless of name.  Using the same
name for two occurrences \emph{constrains} Z3 (both must have the
same interpretation), which can only make queries \emph{easier} to
decide (fewer models to search), never harder.
\end{proposition}

\begin{lstlisting}[caption={Shared function registry: complete implementation}]
class Theory:
    def __init__(self):
        # ... PyObj ADT, TypeTag, RecFunctions ...
        self._shared_fns = {}  # (name, arity) -> Z3 Function

    def shared_fn(self, name, arity):
        """Return canonical Z3 function for a Python operation.

        Both programs use the SAME Z3 function for the same op.
        This is the univalence principle in action.
        """
        key = (name, arity)
        if key not in self._shared_fns:
            sorts = [self.S] * arity + [self.S]
            self._shared_fns[key] = Function(
                f'py_{name}', *sorts)
        return self._shared_fns[key]

    def semantic_axioms(self, solver):
        """Add equational axioms (path constructors)."""
        x = Const('_ax', self.S)
        # sorted: idempotent
        if ('sorted', 1) in self._shared_fns:
            sf = self._shared_fns[('sorted', 1)]
            solver.add(ForAll([x], sf(sf(x)) == sf(x)))
        # reversed: involution
        if ('reversed', 1) in self._shared_fns:
            rf = self._shared_fns[('reversed', 1)]
            solver.add(ForAll([x], rf(rf(x)) == x))
        # len: preserved by sorted, reversed
        if ('len', 1) in self._shared_fns:
            lf = self._shared_fns[('len', 1)]
            for dep in ('sorted', 'reversed', 'list'):
                if (dep, 1) in self._shared_fns:
                    df = self._shared_fns[(dep, 1)]
                    solver.add(ForAll([x],
                        lf(df(x)) == lf(x)))

# Usage in the compiler:
SHARED_BUILTINS = {
    'sorted','reversed','list','tuple','set',
    'frozenset','dict','sum','any','all',
    'enumerate','zip','map','filter',
}

def compile_call(node, env):
    T = env.T
    name = node.func.id
    args = [compile_expr(a, env) for a in node.args]
    if name in SHARED_BUILTINS and args:
        return T.shared_fn(name, len(args))(*args)
    # ... other cases ...
\end{lstlisting}

\begin{lstlisting}[caption={unop: complete Z3 implementation}]
NEG=50; ABS_=52; NOT=53; BOOL_=56; INT_=57; LEN_=55

unop = RecFunction('unop', IntSort(), S, S)
op = Int('_uo'); a = Const('_ua', S)
ai = S.ival(a)

RecAddDefinition(unop, [op, a],
    If(S.is_IntObj(a),
        If(op==NEG,  S.IntObj(-ai),
        If(op==ABS_, S.IntObj(If(ai>=0, ai, -ai)),
        If(op==NOT,  S.BoolObj(ai==0),
        If(op==BOOL_,S.BoolObj(ai!=0),
        If(op==INT_, a,
        S.Bottom))))),
    If(S.is_BoolObj(a),
        If(op==NOT,  S.BoolObj(Not(S.bval(a))),
        If(op==BOOL_,a,
        If(op==INT_, S.IntObj(If(S.bval(a),1,0)),
        S.Bottom))),
    If(S.is_StrObj(a),
        If(op==LEN_, S.IntObj(Length(S.sval(a))),
        If(op==BOOL_,S.BoolObj(Length(S.sval(a))>0),
        S.Bottom)),
    If(S.is_NoneObj(a),
        If(op==BOOL_,S.BoolObj(False),
        If(op==NOT,  S.BoolObj(True),
        S.Bottom)),
    If(S.is_Pair(a),
        If(op==BOOL_,S.BoolObj(True),
        S.Bottom),
    S.Bottom))))))
\end{lstlisting}

\subsection{The Extensibility Principle}

\begin{theorem}[Safe Extension of the Shared Registry]
\label{thm:safe-extension}
Adding a new shared function $\mathsf{py\_new\_op} : \PyObj^k \to
\PyObj$ to the registry $\mathcal{R}$ is \emph{always safe}:
\begin{enumerate}
  \item It does not break decidability of existing queries (the new
    function is uninterpreted; it adds no quantifiers).
  \item It can only \emph{increase} the number of provable
    equivalences (if both programs call the same operation, they
    now get the same Z3 symbol).
  \item Adding axioms for the new function is \emph{optional} ---
    they increase expressiveness but may exit the QF fragment.
\end{enumerate}
\end{theorem}

This means the theory can grow incrementally: every new Python
operation that appears in benchmark programs can be added as a shared
function without auditing the entire theory for decidability regressions.


\chapter{The Semantic Algebra: Complete Axiomatization}
\label{ch:axioms}

This chapter develops the full algebraic theory underlying CTTT.
Each axiom serves a dual role: it is an \emph{equational law} of
Python's operations AND a \emph{path constructor} in the cubical HIT
$\mathsf{PyComp}$ (Definition~\ref{def:pycomp}).  The axioms are
organized by algebraic structure, with proofs of soundness, explicit
Z3 encodings, and analysis of which benchmark pairs each axiom enables.

\section{The Integer Ring}

The integers under $+$ and $\times$ form a commutative ring.  This
is the most fundamental algebraic structure for program equivalence:
whenever two programs compute the same arithmetic expression via
different orderings or groupings, the ring axioms provide the
connecting path.

\subsection{Ring Axioms as Path Constructors}

\begin{definition}[The Integer Ring Paths]
\label{def:int-ring}
The following paths are added to $\mathsf{PyComp}$:

\textbf{Additive abelian group}:
\begin{align}
  \mathsf{add\_comm} &: \prod_{a,b : \PyObj}
    \Path(\mathsf{binop}(\mathsf{ADD}, a, b),\;
          \mathsf{binop}(\mathsf{ADD}, b, a))
    \tag{R1} \\
  \mathsf{add\_assoc} &: \prod_{a,b,c}
    \Path(\mathsf{binop}(\mathsf{ADD}, \mathsf{binop}(\mathsf{ADD}, a, b), c),\;
          \mathsf{binop}(\mathsf{ADD}, a, \mathsf{binop}(\mathsf{ADD}, b, c)))
    \tag{R2} \\
  \mathsf{add\_zero\_r} &: \prod_a
    \Path(\mathsf{binop}(\mathsf{ADD}, a, \mathsf{IntObj}(0)),\; a)
    \tag{R3} \\
  \mathsf{add\_zero\_l} &: \prod_a
    \Path(\mathsf{binop}(\mathsf{ADD}, \mathsf{IntObj}(0), a),\; a)
    \tag{R4} \\
  \mathsf{add\_neg} &: \prod_a
    \Path(\mathsf{binop}(\mathsf{ADD}, a, \mathsf{unop}(\mathsf{NEG}, a)),\;
          \mathsf{IntObj}(0))
    \tag{R5}
\end{align}

\textbf{Multiplicative commutative monoid}:
\begin{align}
  \mathsf{mul\_comm} &: \prod_{a,b}
    \Path(\mathsf{binop}(\mathsf{MUL}, a, b),\;
          \mathsf{binop}(\mathsf{MUL}, b, a))
    \tag{R6} \\
  \mathsf{mul\_assoc} &: \prod_{a,b,c}
    \Path(\mathsf{binop}(\mathsf{MUL}, \mathsf{binop}(\mathsf{MUL}, a, b), c),\;
          \mathsf{binop}(\mathsf{MUL}, a, \mathsf{binop}(\mathsf{MUL}, b, c)))
    \tag{R7} \\
  \mathsf{mul\_one\_r} &: \prod_a
    \Path(\mathsf{binop}(\mathsf{MUL}, a, \mathsf{IntObj}(1)),\; a)
    \tag{R8} \\
  \mathsf{mul\_one\_l} &: \prod_a
    \Path(\mathsf{binop}(\mathsf{MUL}, \mathsf{IntObj}(1), a),\; a)
    \tag{R9} \\
  \mathsf{mul\_zero} &: \prod_a
    \Path(\mathsf{binop}(\mathsf{MUL}, a, \mathsf{IntObj}(0)),\;
          \mathsf{IntObj}(0))
    \tag{R10}
\end{align}

\textbf{Distributivity}:
\begin{align}
  \mathsf{distrib\_l} &: \prod_{a,b,c}
    \Path(\mathsf{binop}(\mathsf{MUL}, a,
      \mathsf{binop}(\mathsf{ADD}, b, c)),\;
      \mathsf{binop}(\mathsf{ADD},
        \mathsf{binop}(\mathsf{MUL}, a, b),
        \mathsf{binop}(\mathsf{MUL}, a, c)))
    \tag{R11}
\end{align}

\textbf{Involution}:
\begin{align}
  \mathsf{neg\_neg} &: \prod_a
    \Path(\mathsf{unop}(\mathsf{NEG}, \mathsf{unop}(\mathsf{NEG}, a)),\; a)
    \tag{R12} \\
  \mathsf{abs\_idem} &: \prod_a
    \Path(\mathsf{unop}(\mathsf{ABS}, \mathsf{unop}(\mathsf{ABS}, a)),\;
          \mathsf{unop}(\mathsf{ABS}, a))
    \tag{R13} \\
  \mathsf{abs\_neg} &: \prod_a
    \Path(\mathsf{unop}(\mathsf{ABS}, \mathsf{unop}(\mathsf{NEG}, a)),\;
          \mathsf{unop}(\mathsf{ABS}, a))
    \tag{R14}
\end{align}
\end{definition}

\begin{theorem}[Soundness of Ring Paths]
Each path in Definition~\ref{def:int-ring} is sound: it holds for
all $\PyObj$ values on the $\mathsf{TInt}$ fiber.
\end{theorem}

\begin{proof}
On the Int fiber, $a = \mathsf{IntObj}(m)$, $b = \mathsf{IntObj}(n)$.
The $\mathsf{binop}$ RecFunction reduces to Z3 native integer
arithmetic: $\mathsf{binop}(\mathsf{ADD}, \mathsf{IntObj}(m),
\mathsf{IntObj}(n)) = \mathsf{IntObj}(m + n)$.  The axioms then
follow from Z3's built-in integer theory (commutativity, associativity,
etc.\ are tautologies of LIA/NIA).
\end{proof}

\subsection{Z3 Encoding}

The ring axioms are NOT added as $\mathsf{ForAll}$ constraints (which
could cause E-matching loops).  Instead, they are applied as
\emph{rewrite rules} during normalization (Phase 2 of
Definition~\ref{def:cnf}):

\begin{lstlisting}[caption={Ring axiom application during normalization}]
def normalize_arithmetic(term, S):
    """Apply ring axioms as rewrite rules."""
    if is_binop(term, ADD):
        a, b = term.arg(1), term.arg(2)
        # R1: canonical ordering (commutativity)
        if term_order(b) < term_order(a):
            term = S.binop(ADD, b, a)
        # R3/R4: identity elimination
        if is_zero(a): return b
        if is_zero(b): return a
    if is_binop(term, MUL):
        a, b = term.arg(1), term.arg(2)
        # R6: canonical ordering
        if term_order(b) < term_order(a):
            term = S.binop(MUL, b, a)
        # R8/R9: identity elimination
        if is_one(a): return b
        if is_one(b): return a
        # R10: zero absorption
        if is_zero(a) or is_zero(b):
            return S.IntObj(IntVal(0))
    # R12: double negation
    if is_unop(term, NEG) and is_unop(term.arg(1), NEG):
        return term.arg(1).arg(1)
    return term
\end{lstlisting}

This is decidable (it's syntactic pattern matching on Z3 terms,
terminating because each rule reduces term complexity).

\subsection{What the Ring Axioms Enable}

The ring axioms are the path constructors for the D1
(Rec$\leftrightarrow$Iter) dichotomy.  They enable:

\begin{itemize}
  \item Factorial: $n \times f(n-1) \equiv f(n-1) \times n$ (R6).
  \item Sum: $n + f(n-1) \equiv f(n-1) + n$ (R1).
  \item Product: $\prod_{i=1}^n a_i$ in any order (R6 + R7 iterated).
  \item Polynomial equivalence: $n(n+1)/2 \equiv 1 + 2 + \ldots + n$
    (R1 + R2 + R11 iterated).
\end{itemize}

\begin{proposition}[Ring Paths for Fold Equivalence]
If the fold operation $\op$ is in $\{+, \times\}$ and the step
functions of two programs differ only by operand order, the ring
commutativity paths (R1 or R6) connect them.

More generally: any two ways of computing a polynomial over the
input variables that give the same polynomial (by the polynomial
identity) are connected by a finite sequence of ring paths.
\end{proposition}


\section{The Lattice of Min and Max}

The operations $\min$ and $\max$ form a \emph{distributive lattice}
on the integers (with the usual ordering).

\begin{definition}[Min-Max Lattice Paths]
\begin{align}
  \mathsf{min\_comm} &: \prod_{a,b}
    \Path(\mathsf{binop}(\mathsf{MIN}, a, b),\;
          \mathsf{binop}(\mathsf{MIN}, b, a))
    \tag{L1} \\
  \mathsf{max\_comm} &: \prod_{a,b}
    \Path(\mathsf{binop}(\mathsf{MAX}, a, b),\;
          \mathsf{binop}(\mathsf{MAX}, b, a))
    \tag{L2} \\
  \mathsf{min\_assoc} &: \prod_{a,b,c}
    \Path(\mathsf{binop}(\mathsf{MIN},
      \mathsf{binop}(\mathsf{MIN}, a, b), c),\;
      \mathsf{binop}(\mathsf{MIN}, a,
      \mathsf{binop}(\mathsf{MIN}, b, c)))
    \tag{L3} \\
  \mathsf{max\_assoc} &: \prod_{a,b,c} \ldots
    \tag{L4} \\
  \mathsf{min\_idem} &: \prod_a
    \Path(\mathsf{binop}(\mathsf{MIN}, a, a),\; a)
    \tag{L5} \\
  \mathsf{max\_idem} &: \prod_a
    \Path(\mathsf{binop}(\mathsf{MAX}, a, a),\; a)
    \tag{L6} \\
  \mathsf{min\_max\_absorb} &: \prod_{a,b}
    \Path(\mathsf{binop}(\mathsf{MIN}, a,
      \mathsf{binop}(\mathsf{MAX}, a, b)),\; a)
    \tag{L7} \\
  \mathsf{max\_min\_absorb} &: \prod_{a,b}
    \Path(\mathsf{binop}(\mathsf{MAX}, a,
      \mathsf{binop}(\mathsf{MIN}, a, b)),\; a)
    \tag{L8} \\
  \mathsf{min\_max\_distrib} &: \prod_{a,b,c}
    \Path(\mathsf{binop}(\mathsf{MIN}, a,
      \mathsf{binop}(\mathsf{MAX}, b, c)),\;
      \mathsf{binop}(\mathsf{MAX},
        \mathsf{binop}(\mathsf{MIN}, a, b),
        \mathsf{binop}(\mathsf{MIN}, a, c)))
    \tag{L9}
\end{align}
\end{definition}

These paths are critical for programs that compute minimums/maximums
in different orders (e.g., DP algorithms that compute
$\min(\min(a, b), c)$ vs.\ $\min(a, \min(b, c))$).

\begin{theorem}[Lattice Paths for DP Equivalence]
\label{thm:lattice-dp}
If two DP algorithms compute the same recurrence using $\min$ or
$\max$ but with different association/ordering of the terms, the
lattice paths (L1--L9) connect them.

In particular: the edit distance recurrence
$d(i,j) = \min(d(i{-}1,j)+1,\, d(i,j{-}1)+1,\, d(i{-}1,j{-}1)+c)$
is invariant under permutation of the three arguments to $\min$
(by L1 + L3 iterated).
\end{theorem}


\section{The Sequence Algebra}

Sequences (lists, tuples) form a \emph{free monoid} under
concatenation, with additional structure from sorting, reversal,
and length.

\subsection{The Free Monoid}

\begin{definition}[Sequence Monoid Paths]
\begin{align}
  \mathsf{concat\_assoc} &: \prod_{x,y,z}
    \Path((x +\!\!+ y) +\!\!+ z,\;
          x +\!\!+ (y +\!\!+ z))
    \tag{S1} \\
  \mathsf{concat\_empty\_r} &: \prod_x
    \Path(x +\!\!+ [],\; x)
    \tag{S2} \\
  \mathsf{concat\_empty\_l} &: \prod_x
    \Path([] +\!\!+ x,\; x)
    \tag{S3}
\end{align}
where $[] = \mathsf{NoneObj}$ (the empty Pair-chain) and
$+\!\!+$ is modeled by the shared function $\mathsf{py\_call\_\_\_add\_\_}$
on Pair-chains.
\end{definition}

\subsection{Sorting as a Retraction}

\begin{definition}[Sorting Paths]
\begin{align}
  \mathsf{sorted\_idem} &: \prod_x
    \Path(\mathsf{py\_sorted}(\mathsf{py\_sorted}(x)),\;
          \mathsf{py\_sorted}(x))
    \tag{S4} \\
  \mathsf{sorted\_perm\_inv} &: \prod_{x,y}
    \Path(\mathsf{py\_sorted}(x +\!\!+ y),\;
          \mathsf{py\_sorted}(y +\!\!+ x))
    \tag{S5} \\
  \mathsf{sorted\_len} &: \prod_x
    \Path(\mathsf{py\_len}(\mathsf{py\_sorted}(x)),\;
          \mathsf{py\_len}(x))
    \tag{S6}
\end{align}
\end{definition}

\begin{remark}[Sorting as a Retraction in CTTT]
In cubical type theory, $\mathsf{sorted}$ is a \emph{retraction}
of the sequence type onto the subtype of sorted sequences:
\[
  r : \mathsf{Seq} \to \mathsf{SortedSeq} \hookrightarrow \mathsf{Seq}
\]
with $r \circ r = r$ (idempotence, path S4).  The retraction forgets
the ordering but preserves the multiset of elements.

This means: in the \emph{quotient type} $\mathsf{Seq}/\sim_\mathsf{perm}$
(sequences modulo permutation), $\mathsf{sorted}$ is the
\emph{canonical representative} of each equivalence class.  Two
sequences with the same sorted form are permutations of each other.

Path S5 says: the sorted form is permutation-invariant.  This is the
key property that makes sorting-based algorithms equivalent to
non-sorting ones when the output only depends on the multiset.
\end{remark}

\subsection{Reversal as an Involution}

\begin{definition}[Reversal Paths]
\begin{align}
  \mathsf{reversed\_invol} &: \prod_x
    \Path(\mathsf{py\_reversed}(\mathsf{py\_reversed}(x)),\; x)
    \tag{S7} \\
  \mathsf{reversed\_len} &: \prod_x
    \Path(\mathsf{py\_len}(\mathsf{py\_reversed}(x)),\;
          \mathsf{py\_len}(x))
    \tag{S8} \\
  \mathsf{sorted\_reversed} &: \prod_x
    \Path(\mathsf{py\_sorted}(\mathsf{py\_reversed}(x)),\;
          \mathsf{py\_sorted}(x))
    \tag{S9}
\end{align}
\end{definition}

S7 makes $\mathsf{reversed}$ a $\Int/2\Int$-action on sequences:
applying it twice returns to the identity.  In cubical type theory,
this is a \emph{loop} in the function space:
$\mathsf{reversed} \cdot \mathsf{reversed} = \mathsf{id}$ (a
2-torsion element of $\pi_1$).

S9 says sorting absorbs reversal: the sorted form of a reversed list
equals the sorted form of the original.  This is because sorting
is permutation-invariant (S5), and reversal is a permutation.

\subsection{Length as a Monoid Homomorphism}

\begin{definition}[Length Paths]
\begin{align}
  \mathsf{len\_additive} &: \prod_{x,y}
    \Path(\mathsf{py\_len}(x +\!\!+ y),\;
          \mathsf{binop}(\mathsf{ADD}, \mathsf{py\_len}(x),
          \mathsf{py\_len}(y)))
    \tag{S10} \\
  \mathsf{len\_empty} &:
    \Path(\mathsf{py\_len}(\mathsf{NoneObj}),\; \mathsf{IntObj}(0))
    \tag{S11} \\
  \mathsf{len\_singleton} &: \prod_a
    \Path(\mathsf{py\_len}(\mathsf{Pair}(a, \mathsf{NoneObj})),\;
          \mathsf{IntObj}(1))
    \tag{S12}
\end{align}
\end{definition}

$\mathsf{len}$ is a \emph{monoid homomorphism} from
$(\mathsf{Seq}, +\!\!+, [])$ to $(\Nat, +, 0)$.  This is a genuine
algebraic structure preservation --- it means Z3 can reason about
lengths using the additive monoid theory (decidable QF\_LIA) even
when the sequences themselves are symbolic.

\subsection{Fold as Universal Catamorphism}

The most powerful algebraic structure: $\fold$ is the
\emph{unique catamorphism} from the free monoid of sequences to
any algebra.

\begin{definition}[Fold Algebra Paths]
\label{def:fold-algebra}
\begin{align}
  \mathsf{fold\_empty} &: \prod_{\op, e}
    \Path(\fold(\op, 0, 0, e),\; e)
    \tag{F1} \\
  \mathsf{fold\_comm} &: \prod_{\op, e, x, y}
    \Path(\fold(\op, 0, 2, e)|_{[x,y]},\;
          \fold(\op, 0, 2, e)|_{[y,x]})
    \tag{F2} \\
  &\quad \text{when $\op$ is commutative} \nonumber \\
  \mathsf{fold\_assoc} &: \prod_{\op, e, x, y, z}
    \Path(\mathsf{binop}(\op, \fold(\op, \ldots)|_{[x,y]}, z),\;
          \fold(\op, \ldots)|_{[x,y,z]})
    \tag{F3} \\
  &\quad \text{when $\op$ is associative} \nonumber \\
  \mathsf{fold\_fusion} &: \prod_{h, \op_1, e_1, \op_2, e_2}
    \Path(h(\fold(\op_1, e_1, \ldots)),\;
          \fold(\op_2, h(e_1), \ldots))
    \tag{F4} \\
  &\quad \text{when $h$ is a homomorphism from $(\op_1, e_1)$ to
    $(\op_2, e_2)$} \nonumber
\end{align}
\end{definition}

\begin{theorem}[Fold Universality]
\label{thm:fold-universal}
For any monoid $(M, \oplus, e)$, the function
$\fold(\oplus, \mathsf{start}, \mathsf{stop}, e)$ is the
\emph{unique} monoid homomorphism from $(\Nat, +, 0)$ to $(M, \oplus, e)$
that maps $1 \mapsto \mathsf{IntObj}(\mathsf{start})$.

This uniqueness is the Nat-recursor path (D1): any two computations
that fold the same operation with the same identity produce the same
result, regardless of fold direction or order (when the operation
is commutative and associative).
\end{theorem}

\begin{proof}
By the universal property of $\Nat$ as the initial object in the
category of monoid homomorphisms from $(\Nat, +, 0)$.  The fold
is the image of the unique morphism.  Any other computation satisfying
the same recurrence factors through this unique morphism, hence
equals the fold.
\end{proof}

\subsection{Sum and Product as Specific Folds}

\begin{corollary}[Sum Permutation Invariance]
$\mathsf{py\_sum}$ is permutation-invariant:
$\mathsf{py\_sum}(\mathsf{py\_sorted}(x)) \equiv \mathsf{py\_sum}(x)$
for all $x$.
\end{corollary}

\begin{proof}
$\mathsf{sum}$ is $\fold(+, 0, \ldots)$.  Since $(+)$ is commutative
and associative (R1, R2), the fold is permutation-invariant (F2, F3).
Sorting is a permutation.  Therefore
$\mathsf{sum}(\mathsf{sorted}(x)) = \mathsf{sum}(x)$.
\end{proof}

Similarly for $\mathsf{product}$, $\mathsf{min}$, $\mathsf{max}$,
$\mathsf{gcd}$, $\mathsf{lcm}$, $\mathsf{xor}$ --- any fold over
a commutative and associative operation.


\section{The Boolean Algebra}

Python's boolean operations form a \emph{Boolean algebra} (actually
a Heyting algebra, since Python 3 doesn't have classical excluded
middle for general predicates, but booleans are classical).

\begin{definition}[Boolean Paths]
\begin{align}
  \mathsf{not\_not} &: \prod_p
    \Path(\mathsf{unop}(\mathsf{NOT}, \mathsf{unop}(\mathsf{NOT}, p)),\; p)
    \tag{B1} \\
  \mathsf{and\_comm} &: \prod_{p,q}
    \Path(\mathsf{BoolObj}(\mathsf{bval}(p) \land \mathsf{bval}(q)),\;
          \mathsf{BoolObj}(\mathsf{bval}(q) \land \mathsf{bval}(p)))
    \tag{B2} \\
  \mathsf{or\_comm} &: \prod_{p,q} \ldots
    \tag{B3} \\
  \mathsf{demorgan\_and} &: \prod_{p,q}
    \Path(\mathsf{unop}(\mathsf{NOT},
      \mathsf{BoolObj}(\mathsf{bval}(p) \land \mathsf{bval}(q))),\;
      \mathsf{BoolObj}(\lnot\mathsf{bval}(p) \lor \lnot\mathsf{bval}(q)))
    \tag{B4} \\
  \mathsf{demorgan\_or} &: \prod_{p,q} \ldots
    \tag{B5} \\
  \mathsf{demorgan\_all} &: \prod_{\mathsf{xs}}
    \Path(\mathsf{py\_all}(\mathsf{xs}),\;
          \mathsf{unop}(\mathsf{NOT},
          \mathsf{py\_any}(\mathsf{py\_map}(\mathsf{NOT}, \mathsf{xs}))))
    \tag{B6} \\
  \mathsf{demorgan\_any} &: \prod_{\mathsf{xs}}
    \Path(\mathsf{py\_any}(\mathsf{xs}),\;
          \mathsf{unop}(\mathsf{NOT},
          \mathsf{py\_all}(\mathsf{py\_map}(\mathsf{NOT}, \mathsf{xs}))))
    \tag{B7}
\end{align}
\end{definition}

B6 and B7 are the \emph{quantifier duality} paths: $\forall x.\, P(x)
\iff \lnot\exists x.\, \lnot P(x)$.  These directly enable the
equivalence of \texttt{all}-based and \texttt{any}-based programs
(benchmark pair eq35).


\section{The Idempotent Semiring}

Several Python operations form \emph{idempotent semirings} --- algebraic
structures where the ``addition'' is idempotent ($a + a = a$).

\begin{definition}[Idempotent Semiring Paths]
For $\min$/$\max$ as ``addition'' and $+$/$\times$ as ``multiplication'':
\begin{align}
  \mathsf{min\_idem} &: \prod_a
    \Path(\mathsf{binop}(\mathsf{MIN}, a, a),\; a)
    \tag{I1} \\
  \mathsf{max\_idem} &: \prod_a
    \Path(\mathsf{binop}(\mathsf{MAX}, a, a),\; a)
    \tag{I2} \\
  \mathsf{set\_idem} &: \prod_x
    \Path(\mathsf{py\_set}(\mathsf{py\_set}(x)),\;
          \mathsf{py\_set}(x))
    \tag{I3} \\
  \mathsf{frozenset\_idem} &: \prod_x
    \Path(\mathsf{py\_frozenset}(\mathsf{py\_frozenset}(x)),\;
          \mathsf{py\_frozenset}(x))
    \tag{I4} \\
  \mathsf{sorted\_idem} &: \prod_x
    \Path(\mathsf{py\_sorted}(\mathsf{py\_sorted}(x)),\;
          \mathsf{py\_sorted}(x))
    \tag{I5} \\
  \mathsf{list\_idem} &: \prod_x
    \Path(\mathsf{py\_list}(\mathsf{py\_list}(x)),\;
          \mathsf{py\_list}(x))
    \tag{I6}
\end{align}
\end{definition}

Idempotence is critical for programs that apply an operation
``just in case'' (e.g., $\mathsf{sorted}(\mathsf{sorted}(x))$ when
one branch already sorts and another doesn't).  The idempotence path
collapses the double application to a single one.

\begin{proposition}[Idempotent Fold]
If $\op$ is idempotent ($\op(a, a) = a$ for all $a$), then:
\[
  \fold(\op, e, \mathsf{xs\_with\_duplicates}) \equiv
  \fold(\op, e, \mathsf{deduplicated\_xs})
\]
Duplicates in the fold input don't affect the result.
\end{proposition}


\section{The Effect Algebra}

The \emph{effect algebra} distinguishes mutation from pure computation.
Unlike the previous algebras (which generate paths), the effect
algebra generates \emph{non-paths}: it asserts that certain terms
are \emph{NOT} equal.

\begin{definition}[Effect Separation Axioms]
\begin{align}
  \mathsf{mut\_ne\_pure} &:
    \mathsf{py\_mut\_iadd}(x, v) \neq
    \mathsf{py\_call\_\_\_add\_\_}(x, v)
    \tag{E1} \\
  \mathsf{mut\_sort\_ne\_sorted} &:
    \mathsf{py\_mut\_sort}(x) \neq \mathsf{py\_sorted}(x)
    \tag{E2} \\
  \mathsf{mut\_reverse\_ne\_reversed} &:
    \mathsf{py\_mut\_reverse}(x) \neq \mathsf{py\_reversed}(x)
    \tag{E3} \\
  \mathsf{ref\_ne\_val} &:
    \mathsf{Ref}(a) \neq \mathsf{IntObj}(n)
    \tag{E4} \\
  \mathsf{ref\_identity} &:
    a \neq b \implies \mathsf{Ref}(a) \neq \mathsf{Ref}(b)
    \tag{E5}
\end{align}
\end{definition}

E1--E3 are the \emph{mutation-purity separation}: in-place operations
and their pure counterparts are \emph{different shared symbols}, so
Z3 treats them as unrelated.  This is how CTTT correctly handles
the 2 FP benchmark cases (neq06, neq20).

E4--E5 are \emph{constructor disjointness}: a Ref is never equal to
a value (different constructors in the ADT), and two Refs with
different addresses are different objects (identity semantics).

\begin{remark}[Effects in Cubical Type Theory]
In standard cubical type theory, the effect axioms are
\emph{negative paths} --- they assert the \emph{absence} of a path.
$\mathsf{mut\_ne\_pure}$ says: there is NO path from
$\mathsf{mut\_iadd}$ to $\mathsf{call\_add}$.

This is formalized using the \emph{empty type} $\mathsf{0}$:
\[
  \mathsf{mut\_ne\_pure} : \Path(\mathsf{mut\_iadd}(x,v),\;
    \mathsf{call\_add}(x,v)) \to \mathsf{0}
\]
--- assuming such a path exists leads to a contradiction.

In Z3, this is implemented trivially: since $\mathsf{mut\_iadd}$ and
$\mathsf{call\_add}$ are different function symbols, Z3 never
identifies them (different uninterpreted functions are always
distinguishable).
\end{remark}


\section{The Composed Algebra}

The full semantic algebra is the \emph{product} of all sub-algebras:

\begin{definition}[The Full Semantic Algebra]
\label{def:full-algebra}
$\mathcal{A}_\mathsf{Py} = \mathcal{R}_\Int \times \mathcal{L}_{\min,\max}
\times \mathcal{M}_\mathsf{Seq} \times \mathcal{B}_\Bool \times
\mathcal{I}_\mathsf{idem} \times \mathcal{E}_\mathsf{eff}$

where:
\begin{center}
\begin{tabular}{l l l l}
  \textbf{Sub-algebra} & \textbf{Paths} & \textbf{Structure} &
  \textbf{Applications} \\
  \hline
  $\mathcal{R}_\Int$ & R1--R14 & Commutative ring &
    Arithmetic, fold order \\
  $\mathcal{L}_{\min,\max}$ & L1--L9 & Distributive lattice &
    DP, optimization \\
  $\mathcal{M}_\mathsf{Seq}$ & S1--S12, F1--F4 &
    Free monoid + sort/rev &
    Collection processing \\
  $\mathcal{B}_\Bool$ & B1--B7 & Boolean algebra &
    Logical equivalence \\
  $\mathcal{I}_\mathsf{idem}$ & I1--I6 & Idempotent semiring &
    Deduplication, caching \\
  $\mathcal{E}_\mathsf{eff}$ & E1--E5 & Effect separation &
    Mutation detection \\
\end{tabular}
\end{center}

Total: \textbf{52 path constructors} (46 positive paths + 5 negative
separation axioms + 1 identity axiom).
\end{definition}

\begin{theorem}[Algebra Soundness]
Every positive path in $\mathcal{A}_\mathsf{Py}$ is sound: it holds
for all $\PyObj$ values on the relevant type fiber.  Every negative
axiom (effect separation) is sound: the separated operations are
genuinely different.
\end{theorem}

\begin{proof}
Each sub-algebra's paths are proved sound in its respective section
above.  The product algebra inherits soundness componentwise (paths
in one sub-algebra don't interact with paths in another, since they
operate on disjoint sets of function symbols).
\end{proof}

\begin{theorem}[Algebra Decidability Preservation]
\label{thm:algebra-decidable}
Adding the paths of $\mathcal{A}_\mathsf{Py}$ to the path search
engine (Definition~\ref{def:cnf}) preserves decidability:
\begin{enumerate}
  \item The rewrite system remains terminating (each rule reduces
    a complexity measure).
  \item The rewrite system remains confluent (by Theorem~\ref{thm:confluence}).
  \item All previously decidable queries remain decidable.
  \item The paths are applied as rewrite rules (no quantifiers
    added to the Z3 solver).
\end{enumerate}
\end{theorem}

\section{Interaction Between Sub-Algebras}

The sub-algebras are not fully independent --- they interact through
the $\PyObj$ ADT.  The interactions provide \emph{additional} paths
beyond what each sub-algebra provides alone:

\subsection{Ring--Fold Interaction}

The fold catamorphism (F1--F4) interacts with the ring axioms (R1--R14)
to produce the \emph{fold canonicalization} path: any fold of a ring
operation produces a canonical form.

\begin{theorem}[Fold Canonicalization]
\label{thm:fold-canon}
For a ring operation $\op \in \{+, \times\}$:
\[
  \fold(\op, \mathsf{start}, \mathsf{stop}, e) =
  \fold(\op, \mathsf{start}', \mathsf{stop}', e')
\]
if and only if they compute the same ring element
(same value as a polynomial in $\mathsf{start}$ and $\mathsf{stop}$).

This is decidable: polynomial identity testing over $\Int$ is in P
(by the Schwartz--Zippel lemma applied to integer polynomials, or
by direct Z3 integer arithmetic).
\end{theorem}

\subsection{Lattice--Fold Interaction}

Folding a lattice operation ($\min$ or $\max$) produces the
lattice join/meet of the elements.  Combined with the lattice
distributivity (L9), this enables:

\begin{proposition}
$\min(\fold(+, \ldots), \fold(+, \ldots))$ can be distributed
over the folds (by L9 + F3), enabling DP-style equivalences
where the ``min'' is applied at different levels of the recursion.
\end{proposition}

\subsection{Sequence--Boolean Interaction}

The De Morgan paths (B6, B7) interact with the sequence fold
structure.  $\mathsf{all}(\mathsf{xs})$ is semantically
$\fold(\land, \top, \mathsf{xs})$ (fold of conjunction), and
$\mathsf{any}(\mathsf{xs})$ is $\fold(\lor, \bot, \mathsf{xs})$
(fold of disjunction).  The De Morgan path transposes these:
\[
  \fold(\land, \top, \mathsf{xs}) \equiv
  \lnot\,\fold(\lor, \bot, \mathsf{map}(\lnot, \mathsf{xs}))
\]

This is a \emph{natural transformation} between two functors on
the sequence category --- a genuinely categorical path that the ring
and lattice algebras alone cannot produce.


